% ID: SEL-D0083-OA-E04
% ESTADO: cotejada. Segunda pasada visual y borrador compilado revisados.
% Fuente DOC: Selectividad/sel_2012_matematicas_aplicadas/5 MATEMÁT. APLIC. CC. SS/EXAMEN 3-ANDALUCÍA.doc
% SHA-256 original: c8688a19515dfeb348e8f8817d3f4a6c0a967723c270f2ec42e7f365d001e555
% Página convertida: 1; opción A; ejercicio 4.

En una caja de ahorros se sabe que el porcentaje de los nuevos clientes que contratan un plan de pensiones no supera el $23\%$. El director de una de las sucursales decide hacer un regalo a cualquier nuevo cliente que contrate uno de esos planes y, tras un mes, comprueba que 110 de los 470 nuevos clientes han contratado un plan de pensiones.
\begin{enumerate}[label=\alph*)]
\item \textbf{(1.5 puntos)} Plantee un contraste de hipótesis, con $H_0:p\leq0.23$, para decidir si, con los datos dados, se puede afirmar que la medida del director ha aumentado la contratación de estos planes de pensiones. Halle la región de aceptación de este contraste de hipótesis para un nivel de significación del $5\%$.
\item \textbf{(1 punto)} Según el resultado del apartado anterior, \textquestiondown{}qué conclusión podemos obtener sobre la medida tomada por el director de esta sucursal?
\end{enumerate}
